
\section{Method}

\subsection{Overview}

MinimalNOVA replaces the four attention projections of each transformer layer 
with \textbf{SVD-factorized linear modules}:

\[
W \to (U, V), \qquad
y = (x V^\top) U^\top + \text{LoRA}(x) + b
\tag{4.1}
\]

The factorization is done independently per projection. Rank is chosen 
adaptively per layer based on spectral energy.

\textbf{Figure 4.1} shows the pipeline.

\begin{lstlisting}
       +------------------------------------------+
       |  Input weight W $\in$ R^{d_out $$\times$ d_in}       |
       +----------------+-------------------------+
                        |
                         
       +------------------------------------------+
       |  Compute SVD: W = U $\Sigma$ V^T                |
       +----------------+-------------------------+
                        |
                         
       +------------------------------------------+
       |  Compute cumulative energy E(r)          |
       |  r_eff = min(r_max, argmin E(r) $\ge$  )     |
       +----------------+-------------------------+
                        |
                         
       +------------------------------------------+
       |  Truncate to r_eff, symmetrize $\Sigma$         |
       |  U_r = U[:, :r]   $\sqrt{}$$\Sigma$_r                   |
       |  V_r = $\sqrt{}$$\Sigma$_r   V[:r, :]^T                 |
       +----------------+-------------------------+
                        |
                         
       +------------------------------------------+
       |  Replace Linear with FactorisedLinear    |
       |  Freeze U, V; add trainable LoRA         |
       +------------------------------------------+
\end{lstlisting}

\subsection{The `FactorisedLinear` Module}

\subsubsection{Forward Pass}

The module computes:

\[
y = \text{up}(\text{down}(x)) + \text{LoRA}(x) + b
\tag{4.2}
\]

where:

\begin{itemize}
\item \texttt{down: nn.Linear(d\_in, r)} -- down-projection, weights $V_r \in \mathbb{R}^{r \times d_\text{in}}$
\item \texttt{up: nn.Linear(r, d\_out)} -- up-projection, weights $U_r \in \mathbb{R}^{d_\text{out} \times r}$
\item \texttt{LoRA(x) = (x A\^{}\top) B\^{}\top \cdot s} -- optional adapter

\end{itemize}
\subsubsection{Parameter Count}

\begin{table}[H]
\centering
\small
\begin{tabular}{ll}
\toprule
\textbf{Configuration} & \textbf{Parameters} \\
\midrule
Original \texttt{nn.Linear} & $d_\text{in} d_\text{out}$ \\
\texttt{FactorisedLinear} at rank $r$ & $r(d_\text{in} + d_\text{out})$ \\
\texttt{FactorisedLinear} + LoRA($r_L$) & $r(d_\text{in} + d_\text{out}) + r_L(d_\text{in} + d_\text{out})$ \\
\bottomrule
\end{tabular}
\end{table}

\textbf{Example} for a $768 \times 768$ attention projection:

\begin{table}[H]
\centering
\small
\begin{tabular}{lll}
\toprule
\textbf{Config} & \textbf{Params} & \textbf{Ratio} \\
\midrule
Original & 589,824 & 1.00$$\times$ \\
r=64 & 98,304 & \textbf{6.00$$\times$} \\
r=128 & 196,608 & 3.00$$\times$ \\
r=64 + LoRA r=16 & 122,880 & 4.80$$\times$ \\
\bottomrule
\end{tabular}
\end{table}

\subsection{Energy-Adaptive Rank Selection}

\subsubsection{Algorithm}

Given weight $W$ and hyperparameters $r_\text{max}$ (cap) and $\eta$ (energy):

\begin{enumerate}
\item Compute SVD: $W = U \Sigma V^\top$.
\item Compute cumulative energy: $E(r) = \sum_{i \leq r} \sigma_i^2 / \sum_i \sigma_i^2$.
\item Find the smallest $r$ with $E(r) \geq \eta$:
  $$
  r\_\text{energy} = \left|\left\{i : E(i) < \eta\right\}\right| + 1
  \tag{4.3}
  $$
\item Cap the rank:
  $$
  r\_\text{eff} = \min(r\_\text{max}, r\_\text{energy}, \min(d\_\text{in}, d\_\text{out}))
  \tag{4.4}
  $$

\end{enumerate}
\subsubsection{Why Energy?}

The energy criterion captures \textbf{how concentrated the spectrum is}. 
Consider two extreme cases:

\begin{itemize}
\item \textbf{Rank-1 matrix} ($\sigma_1 > 0$, $\sigma_{i>1} = 0$): $r_\text{energy} = 1$
  for any $\eta > 0$. Very compressible.
\item \textbf{Identity matrix} ($\sigma_i = 1$ for all $i$): $E(r) = r/d$, so
  $r_\text{energy} = \lceil \eta d \rceil$. Less compressible.

\end{itemize}
Real transformer weights lie between these extremes, with power-law or 
exponential decay. The energy criterion adapts to each matrix's specific 
spectrum.

\subsubsection{Comparison to Fixed Rank}

\begin{table}[H]
\centering
\small
\begin{tabular}{lll}
\toprule
\textbf{Criterion} & \textbf{Pros} & \textbf{Cons} \\
\midrule
Fixed rank & Simple & Ignores heterogeneity \\
Energy & Adaptive & No calibration \\
Fisher-weighted & Activation-aware & Needs data \\
\bottomrule
\end{tabular}
\end{table}

MinimalNOVA uses energy, which is \textbf{training-free} and requires no 
calibration data.

\subsection{Weight Factorization}

\subsubsection{Symmetric Distribution}

We distribute $\Sigma^{1/2}$ between the two factors:

\[
U_r = U_{:, :r} \cdot \Sigma_{:r, :r}^{1/2}, \qquad
V_r = \Sigma_{:r, :r}^{1/2} \cdot V_{:r, :}^\top
\tag{4.5}
\]

This gives $W_r = U_r V_r$ with balanced factor magnitudes, which improves 
training stability.

\subsubsection{Storage in Linear Modules}

The framework stores $V_r$ in \texttt{down.weight} and $U_r$ in \texttt{up.weight}:

\begin{lstlisting}
self.down.weight.data[:r, :] = V_r       # (r, d_in)
self.up.weight.data[:, :r] = U_r         # (d_out, r)
\end{lstlisting}

For $r < \text{target rank}$, the remaining rows/columns are zero.

\subsection{Weight Transfer from Pretrained Attention}

MinimalNOVA extracts $W_Q, W_K, W_V, W_O$ from the pretrained attention 
module. This is nontrivial because different architectures store them 
differently.

\subsubsection{Extraction Path (per architecture)}

\begin{table}[H]
\centering
\small
\begin{tabular}{lllll}
\toprule
\textbf{Architecture} & \textbf{$W_Q$} & \textbf{$W_K$} & \textbf{$W_V$} & \textbf{$W_O$} \\
\midrule
DistilBERT & \texttt{q\_lin} & \texttt{k\_lin} & \texttt{v\_lin} & \texttt{out\_lin} \\
RoBERTa & \texttt{self.query} & \texttt{self.key} & \texttt{self.value} & \texttt{output.dense} \\
GPT-2 & \texttt{c\_attn[:d]} & \texttt{c\_attn[d:2d]} & \texttt{c\_attn[2d:]} & \texttt{c\_proj} \\
Qwen3.5 & \texttt{q\_proj} & \texttt{k\_proj} & \texttt{v\_proj} & \texttt{o\_proj} \\
\bottomrule
\end{tabular}
\end{table}

\subsubsection{GPT-2 Fused c\_attn}

GPT-2 fuses $Q, K, V$ into one \texttt{c\_attn}:

\[
W_{c\_attn} = \begin{pmatrix} W_Q \\ W_K \\ W_V \end{pmatrix} \in \mathbb{R}^{3d \times d}
\tag{4.6}
\]

(Conv1D stores this transposed.) MinimalNOVA splits:

\begin{lstlisting}
w = src.c_attn.weight.t()   # (3d, d)
q_w, k_w, v_w = w.chunk(3, dim=0)
\end{lstlisting}

\subsubsection{RoBERTa Nested}

The output projection $W_O$ is at \texttt{attention.output.dense}, \textbf{outside} the 
inner attention:

\begin{lstlisting}
# Extract Q, K, V from inner self-attention
q_w = src.self.query.weight
k_w = src.self.key.weight
v_w = src.self.value.weight

# Extract W_O from outer layer
w_o = src.output.dense.weight
\end{lstlisting}

When replacing the inner attention, the outer \texttt{output.dense} remains in 
the model and applies the output projection \textbf{after} our module returns. 
So we set \texttt{is\_nested\_encoder = True} and skip \texttt{out\_proj} in the forward 
pass.

\subsubsection{GQA (Qwen)}

For GQA, $H_{kv} < H$. The K and V projections have smaller output 
dimensions:

\[
W_K \in \mathbb{R}^{H_{kv} d_h \times d}, \quad
W_V \in \mathbb{R}^{H_{kv} d_h \times d}
\tag{4.7}
\]

MinimalNOVA factorizes each at its natural dimension. During forward, 
\texttt{repeat\_kv} expands K and V to match H query heads.

\subsection{Validation and Fallback}

\subsubsection{Frobenius Validation}

After factorization, MinimalNOVA does not explicitly validate (unlike A3). 
Instead, it relies on \textbf{weight copy verification}:

\begin{lstlisting}
for proj in ["q", "k", "v"]:
    w = extracted[proj][0]
    if w is not None:
        norm = torch.linalg.norm(w.float())
        if norm < 1e-4:
            raise RuntimeError(f"Near-zero Frobenius norm on {proj}")
\end{lstlisting}

This catches catastrophic extraction failures but does not bound the 
\textbf{factorization error}. This is a limitation ( 10.7).

\subsubsection{Fallback}

If \texttt{copy\_weights} raises, \texttt{inject\_minimal\_nova} catches and returns \texttt{None}. 
The original attention is preserved.

\subsection{LoRA Recovery}

After factorization, base weights are frozen:

\begin{lstlisting}
def freeze_base_weights(self):
    self.down.weight.requires_grad = False
    self.up.weight.requires_grad = False
    if self.up.bias is not None:
        self.up.bias.requires_grad = False
    if self.lora_rank > 0:
        self.lora_A.requires_grad = True
        self.lora_B.requires_grad = True
\end{lstlisting}

Only \texttt{lora\_A} and \texttt{lora\_B} are trained. LoRA $B$ is initialized to zero, so:

\[
y_{\text{init}} = (x V^\top) U^\top + b
\tag{4.8}
\]

which equals the SVD approximation.

\subsection{The Full Algorithm}

\textbf{Algorithm 1:} MinimalNOVA injection
\begin{lstlisting}
Input:  model M, rank_cap r_max, energy  
Output: model M' with factorized attention

1:  for layer   in transformer_layers(M) do
2:      src $\leftarrow$ get_attention( )
3:      extract Q, K, V, O from src  (architecture-dependent)
4:
5:      for (W, target) in [(Q, q_proj), (K, k_proj), (V, v_proj), (O, out_proj)] do
6:          U, $\Sigma$, V $\leftarrow$ svd(W)
7:          E(r) $\leftarrow$ cumsum($\Sigma$ ) / sum($\Sigma$ )
8:          r_eff $\leftarrow$ min(r_max, |{i: E(i) <  }| + 1)
9:          U_r $\leftarrow$ U[:, :r]   $\sqrt{}$$\Sigma$[:r]
10:         V_r $\leftarrow$ $\sqrt{}$$\Sigma$[:r]   V[:r, :]
11:         target.down.weight[:r_eff, :] $\leftarrow$ V_r
12:         target.up.weight[:, :r_eff] $\leftarrow$ U_r
13:     end for
14:
15:     freeze_all_base_factors( )
16:     replace_attention( , new)
17: end for
18: return M'
\end{lstlisting}

\subsection{Complexity Analysis}

\begin{table}[H]
\centering
\small
\begin{tabular}{lll}
\toprule
\textbf{Phase} & \textbf{Time} & \textbf{Memory} \\
\midrule
Weight extraction & $\mathcal{O}(N_l d^2)$ & $\mathcal{O}(d^2)$ \\
SVD per projection & $\mathcal{O}(d^3)$ & $\mathcal{O}(d^2)$ \\
Truncation & $\mathcal{O}(N_l r d^2)$ & $\mathcal{O}(d^2)$ \\
Total & $\mathcal{O}(N_l d^3)$ & $\mathcal{O}(d^2)$ \\
\bottomrule
\end{tabular}
\end{table}

For DistilBERT ($N_l = 6$, $d = 768$): ~10 seconds on RTX 4070.

\subsection{Theoretical Notes}

\subsubsection{Optimality per Rank}

Given rank $r$, the truncated SVD is optimal in Frobenius norm by Eckart--Young.

\subsubsection{Energy Criterion}

\textbf{Proposition.} The energy-adaptive rank satisfies:

\[
\frac{\|W - U_r V_r\|_F^2}{\|W\|_F^2} = 1 - E(r)
\tag{4.9}
\]

\textbf{Proof.} $\|W - W_r\|_F^2 = \sum_{i > r} \sigma_i^2 = \|W\|_F^2 - \sum_{i \leq r} \sigma_i^2$. 
Dividing by $\|W\|_F^2$ gives (4.9). $\square$

This means: choosing $r$ with $E(r) \geq \eta$ guarantees relative 
reconstruction error $\leq \sqrt{1 - \eta}$.

For $\eta = 0.95$: error $\leq \sqrt{0.05} \approx 0.224$.
For $\eta = 0.99$: error $\leq \sqrt{0.01} = 0.10$.

\subsubsection{Why Symmetric Distribution Helps}

Distributing $\Sigma^{1/2}$ symmetrically ensures:
\begin{itemize}
\item Both factors have unit spectral norm (up to constants)
\item Gradient magnitudes are balanced during LoRA training
\item Numerical stability in FP16 (no factor dominates)

\end{itemize}
\subsection{Summary}

MinimalNOVA:
\begin{enumerate}
\item Factorizes each attention projection $W \to (U, V)$.
\item Selects rank via spectral energy $\eta$.
\item Uses LoRA for recovery.
\item Handles four attention architectures.
\item Achieves \textbf{real parameter reduction} of $d^2 / (2rd)$ per projection.
\end{enumerate}
\begin{lstlisting}

---
\end{lstlisting}